% compendium/portfolio-construction-risk/portfolio-construction-models.tex
% Entries migrated on 2026-10-15 from the legacy portfolio-construction-models file (slice 28, sysinv/portconstmod.tex),
% owned by portfolio-construction, the file's single map candidate. Dispositions and source checks in
% notes/migration-slices.md, the file verbatim in notes/artefacts/portfolio-construction-risk/portfolio-construction-models.tex.

\section{Portfolio construction models: rules and optimisers}

\begin{qfmethod}[Rule-based position sizing: equal position, equal risk, alpha-driven and decision-tree weighting, with Narang's case against unequal weights]\label{meth:rule-based-position-sizing}
\qffield{Statement} Narang divides quantitative portfolio construction models into two families: \emph{rule-based} models, heuristics ``derived from human experience, such as by trial and error'', which can be ``exceedingly simple or rather complex''; and \emph{optimisers}, which search for the best value of an objective function \citep[chap.~6, ``Portfolio Construction Models'', pp.~99--100]{s7-narang} (the second family is in \hyperref[mod:mean-variance-portfolio-optimisation]{mean--variance optimisation}). He names four rule-based types (sec.~6.1, p.~100):
\begin{itemize}
\item \emph{equal position weighting}: every position the alpha model deems worth holding gets the same size; signal strength is ignored ``except insofar as the signal is strong enough to be worthy of a position at all'', and the instruments are implicitly assumed homogeneous in risk (sec.~6.1.1, pp.~100--101);
\item \emph{equal risk weighting}: sizes ``inversely to their volatilities (or whatever other measure of risk, such as drawdown, is preferred)'', so that each position is equalised ``by the amount of risk that the allocation contributes to the portfolio'' (sec.~6.1.2, p.~102);
\item \emph{alpha-driven weighting}: size grows with the forecast, under limits from the risk model on single positions and groups, for example 3 percent a position and 20 percent a sector (sec.~6.1.3, pp.~103--104);
\item \emph{decision-tree weighting}: ``a defined set of rules in a particular order to determine position sizing'', named on p.~100 only; the third edition has no section on it.
\end{itemize}

\qffield{Conditions} \emph{Why some quants weight equally.} Narang gives two problems with \emph{unequal} weighting, which the equal-weighters avoid (pp.~101--102). First, it assumes enough statistical power to forecast not only direction but also the relative size or probability of moves. Second, it leads to ``a few large bets on the `best' forecasts and many smaller bets on the less dramatic forecasts'', and so to excess idiosyncratic risk: the strongest momentum or reversal signals often come where prices have already moved most, sometimes on real information (his ``adverse selection bias''), and large forecasts can come from bad data, such as a price quoted in pence instead of pounds. Winsorising forecasts can partly address this; equal weights need no such adjustment. \emph{Equal risk weighting} relies on a backward-looking risk measure: his example is bank stocks, stable for a decade and then more volatile than many technology stocks in 2008 (p.~103). \emph{Alpha-driven weighting} relies on the alpha model being right about size as well as direction; trend followers using it ``can suffer sharp drawdowns relatively frequently'', holding the largest position as the trend reverses (p.~104). Costs can enter any rule, for example by setting a forecast to zero when it is below its expected transaction cost (sec.~6.1.4, p.~104; \hyperref[mod:portfolio-transaction-cost-function]{cost functions}). Liquidity can cap any weight (p.~102).

\qffield{Practice} Interpretation, checked numerically: his Exhibit~6.1 (p.~102, page image) gives volatilities of 2.5 and 2.0 percent and ``volatility-adjusted'' weights of 39 and 61 percent. Those are inverse-\emph{variance} weights, $(1/2.5^2)/(1/2.5^2+1/2.0^2)=39.0\%$; inverse-volatility weights would be 44.4 and 55.6 percent, which give each stock the same $w_i\sigma_i$ (1.11) while the table's weights do not (0.98 and 1.22). Neither scheme in general equalises the contribution to \emph{portfolio} risk, $w_i(\Sigma w)_i/w^\top\Sigma w$, which depends on correlations. For two assets inverse-volatility weights do equalise it at any correlation, but not for more. With three assets of 20 percent volatility, where the first two correlate at 0.8 and the third is uncorrelated, equal (inverse-volatility) weights give risk contributions of 39, 39 and 22 percent. Equal contributions need weights of 29.9, 29.9 and 40.2 percent (checked). Narang's ``contributes to the portfolio'' therefore describes equal-risk-contribution weighting, which uses the full covariance matrix, more than the inverse-volatility rule he gives (stub \texttt{meth:equal-risk-contribution}).

\qffield{Pitfalls} The legacy block put Narang's second argument \emph{against} unequal weighting, ``few large bets on `best' forecast, many smaller bets'', under equal weighting, where it reads as a property of the equal-weight scheme. It describes what equal weighting avoids. The legacy decision-tree details (``constraints may include percentage limits'', ``model size grows dramatically'') are not in the third edition's chapter; they may come from an earlier edition (not checked). Narang's choice between rules is his observation: rule-based users are mostly futures traders with ``intrinsic'' forecasts, and optimiser users are mostly equity market-neutral traders with relative forecasts (sec.~6.4, p.~123). He gives no data for this.

\qffield{See also} \hyperref[mod:mean-variance-portfolio-optimisation]{mean--variance optimisation} (the other family); \hyperref[meth:signal-portfolio-allocation]{signal-portfolio allocation} (rule-based signal portfolios mixed by an optimiser); \hyperref[mod:portfolio-transaction-cost-function]{cost functions} (costs as inputs); \hyperref[def:common-investor-contagion]{common-investor contagion} (risk targeting under rising volatility).

\qffield{Further reading} Equal risk contribution and volatility targeting belong to the risk-parity pathway, whose readings are unread; a stub here names it. The legacy file it replaces is kept verbatim in the section's folder of \nolinkurl{notes/artefacts/}, as\\ \nolinkurl{portfolio-construction-models.tex}.

\qffield{Index} portfolio construction model; rule-based weighting; equal weighting; equal position weighting; equal risk weighting; inverse volatility; inverse variance; risk contribution; equal risk contribution; alpha-driven weighting; position limits; sector limits; decision tree; winsorisation; adverse selection; bad data; intrinsic and relative alpha

\qffield{Status} Verification: partly checked 2026-10-15 (against the text layer of a local copy of \texttt{s7-narang}, third edition, Wiley 2024, chap.~6, pp.~99--104 and~123, with Exhibit~6.1 on p.~102 by page image; the weights and risk contributions computed). Stability: draft. Currency: practitioner account as of 2024; the risk-contribution arithmetic is timeless.
\end{qfmethod}

\begin{qfmodel}[Mean--variance portfolio optimisation: objectives, inputs, constraints and solvers, and what the legacy claimed about them]\label{mod:mean-variance-portfolio-optimisation}
\qffield{Statement} Markowitz's problem, as written by Boyd, Johansson, Kahn, Schiele and Schmelzer, chooses the weights $w\in\mathbb R^n$ given expected returns $\mu$ and a covariance matrix $\Sigma$ in one of two forms \citep[author version of 5 Jan.\ 2024, sec.~1.1, eqs.~(1)--(2), pp.~4--5]{wide-boyd-markowitz-seventy}:
\[ \max_w\ \mu^\top w\ \ \text{s.t.}\ \ w^\top\Sigma w\le\sigma_{\mathrm{tar}}^2,\ \mathbf 1^\top w=1;\qquad \max_w\ \mu^\top w-\gamma\,w^\top\Sigma w\ \ \text{s.t.}\ \ \mathbf 1^\top w=1. \]
Both trace the same trade-off curve of Pareto-optimal portfolios as $\sigma_{\mathrm{tar}}$ or $\gamma>0$ varies, and both have closed forms; that of the second is $w^\star=\tfrac1{2\gamma}\Sigma^{-1}(\mu+\nu^\star\mathbf 1)$, whose $\Sigma^{-1}$ hints at sensitivity to the inputs. Adding $w\ge0$ (long only) makes the second a quadratic programme (pp.~4--5). Maximising the Sharpe ratio is a third objective, also solvable by convex optimisation ``at least when the constraints are simple'' (sec.~3.2, p.~20). Narang describes the same model in words: inputs of expected returns, volatilities and a correlation (covariance) matrix; output the efficient frontier; a common objective of ``the highest possible return for a unit of risk''; and in practice also the portfolio's size in currency, a desired risk level ``usually measured in terms of volatility or expected drawdown'', and other constraints such as a hard-to-borrow list, which ``are not required by the optimizer'' \citep[chap.~6, sec.~6.2, pp.~104--106]{s7-narang}.

\qffield{Notation} $\mu$, $\Sigma$, $w$, $\gamma$ (risk aversion), $\sigma_{\mathrm{tar}}$ (target volatility) as in Boyd et al.; the Sharpe ratio $\mu^\top w/\sqrt{w^\top\Sigma w}$ (excess returns, no benchmark) is unchanged when $w$ is scaled by any positive number.

\qffield{Conditions} \emph{Which objective.} Mean--variance maximises return for a variance target, or return minus a variance penalty; maximising the Sharpe ratio selects one direction on the frontier (with cash and leverage, the tangency portfolio, Narang's Exhibit~6.7, p.~122, page image). These are different problems with different solutions. Expected drawdown is not part of the mean--variance model: Narang lists it among practical inputs that set a risk level (p.~105) and names a return-to-drawdown objective as a different choice (p.~124). The legacy ``desired risk level (volatility or expected drawdown)'' and ``maximising return of portfolio relative to volatility'' put both inside one model. \emph{Unconstrained optimisation.} Narang calls it ``a quirk of unconstrained optimizers'' that they ``often'' put everything in the instrument with the highest risk-adjusted return (sec.~6.2.2.1, p.~114); the legacy made this ``might''. It is not a property of mean--variance: the unconstrained solution is proportional to $\Sigma^{-1}$ terms and generally holds every asset, and his own unconstrained two-stock example has its maximum Sharpe ratio at about 85 percent CVX and 75 percent NFLX (p.~106; Exhibit~6.2, p.~107, page image). Interpretation: a single-asset answer arises when only expected return is maximised under a budget and no shorting, or when risk is ignored. Because the Sharpe ratio is scale-free, the point he marks is one of a ray of equally good portfolios within his $\pm100$ percent search box (checked: $\Sigma^{-1}\mu$ and three times it have the same ratio).

\qffield{Practice} \emph{Inputs.} Expected returns come from the alpha model; a sign-only forecast can be used by setting all positive and all negative forecasts equal (p.~108). For volatility Narang calls GARCH, with Bollerslev's 1986 paper, ``the most widely used such technique (particularly in equity trading)'', and adds that a simpler EWMA ``may perform just about as well'' (pp.~108--110); he cites no evidence for the prevalence. Boyd et al.\ list EWMA, dynamic conditional correlation and iterated EWMA as covariance estimators that work well in practice (sec.~3.4, p.~22). For correlations Narang shows instability: S\&P~500 against Nikkei~225 rolling one-year correlations from $+0.01$ to $+0.66$ (pp.~110--111, Exhibit~6.4 by page image), and XOM--CVX steadier than XOM--NFLX (Exhibit~6.5, p.~112), and says grouping similar instruments, ``in the definition of \emph{relative} in a relative alpha strategy and/or in the specification of the risk model'', improves stability (p.~111). \emph{Constraints} limit positions and groups, and can end up driving the portfolio: with 100 names and a 1.5 percent cap the best position is only 1.5 times the average (sec.~6.2.2.2, p.~114). Risk penalties can rise faster than linearly, and costs enter as per-stock or generic impact functions (pp.~114--115); Boyd et al.\ use a spread term plus a $\lvert z\rvert^{3/2}$ impact term (sec.~2.4, eq.~(10), p.~17), the form in \hyperref[mod:portfolio-transaction-cost-function]{cost functions}. Industry-neutral constraints can turn a positive forecast into a short, and the substitution effect can prefer a cheaper correlated stock (sec.~6.2.3, pp.~120--121). The output is a target portfolio; the trades are its difference from the current one (sec.~6.3, pp.~122--123). \emph{Solvers.} Convex formulations, including nondifferentiable terms, are solved reliably and fast; integer shares, cardinality limits and minimum position sizes make the problem nonconvex and much slower (Boyd et al., secs.~1.4 and~4.3, pp.~9--10 and 27--28). Narang's ``10{,}000 assets \dots\ in about one second'' (p.~120) matches Boyd et al.'s Table~2 only with a factor model: 0.65--0.89 seconds for 10{,}000 assets and 50--100 factors with one commercial solver (sec.~5.6, p.~40).

\qffield{Pitfalls} \emph{Estimation error.} Boyd et al.\ agree that the basic method is sensitive to its inputs, and answer with regularisation and robust optimisation: long-only or leverage limits, Black--Litterman's shrinkage of the mean towards market-implied returns, and covariance shrinkage (sec.~1.3, pp.~8--9; sec.~3.4, p.~22). \emph{Black--Litterman.} Narang credits it with a default portfolio and a blend of ``an investor's expectations with a degree of confidence about those expectations with the historical precedent evident in the data'' (sec.~6.2.2.3, p.~117). In Boyd et al.'s account it acts on the expected returns, towards the market-implied return, not on correlations. The legacy ``adjusts historically observed correlation levels by utilising investor's forecast'' is in neither source and is not adopted. Black and Litterman's paper is unread (stub); Narang's note~5 dates it 1982, his text and the catalogue 1992. \emph{Resampled efficiency.} Narang presents Michaud's method as reordering the observed returns by Monte Carlo ``to reduce the estimation error inherent in the inputs'' (sec.~6.2.2.5, pp.~118--119). His own remark that average return and volatility ``remain the same across all these alternate histories'' shows that reordering cannot change the sample means and covariances an optimiser uses. His example concerns the robustness of a trading strategy, not optimiser inputs. What Michaud's procedure does and whether it improves out-of-sample results were not checked (stub). \emph{Machine learning.} The legacy ``data-mining approaches: supervised learning or genetic algorithms, as MVO involves searching many possible portfolios'' is not in the third edition. There Narang says ``we do not see many strictly `machine learning' models for portfolio construction'' and that the state of the art remains Markowitz-type optimisation (sec.~6.2.2, pp.~113--114). Interpretation: supervised learning fits a forecast; it does not choose weights. A convex portfolio problem has solvers that reach its optimum, so a generic search heuristic such as a genetic algorithm is no justification there; such heuristics belong to nonconvex problems. Not adopted. \emph{Narang's details not adopted.} He describes convex functions as ``smooth and continuous'' with ``a single highest or lowest point'' (p.~120); convex functions can have kinks, such as $\lvert x\rvert$ or the costs above, and flat sets of minima. He calls GARCH a stochastic-volatility model and says a negative autoregression ``implies \dots\ trending behavior'' (pp.~108--110); neither is adopted. His 1{,}998{,}000 distinct correlations for 2{,}000 instruments (p.~113) should be $2000\times1999/2=1{,}999{,}000$ (checked). His grouping evidence is two pairs over one sample, not a test of a method; factor and shrinkage estimators are stubbed.

\qffield{See also} \hyperref[meth:rule-based-position-sizing]{rule-based position sizing}; \hyperref[meth:signal-portfolio-allocation]{signal-portfolio allocation}; \hyperref[mod:portfolio-transaction-cost-function]{cost functions}; \hyperref[mod:almgren-chriss-execution]{mean--variance execution} (the same trade-off over an execution schedule); \hyperref[meth:backtest-evaluation-diagnostics]{backtest evaluation diagnostics} (Sharpe and Calmar ratios as Narang defines them).

\qffield{Further reading} Markowitz's paper, Black and Litterman, Michaud, shrinkage and factor-model readings of the pathway: not read (stubs). The legacy file it replaces is kept verbatim in the section's folder of \nolinkurl{notes/artefacts/}, as\\ \nolinkurl{portfolio-construction-models.tex}.

\qffield{Index} mean--variance optimisation; Markowitz; modern portfolio theory; efficient frontier; risk aversion; target volatility; Sharpe ratio maximisation; tangency portfolio; unconstrained optimisation; constrained optimisation; position limits; industry neutrality; substitution effect; expected returns; sign forecasts; GARCH; EWMA; correlation instability; estimation error; regularisation; robust optimisation; Black--Litterman; resampled efficiency; convex optimisation; mixed-integer; solve time; target portfolio

\qffield{Status} Verification: partly checked 2026-10-15 (against the text layer of a local copy of \texttt{s7-narang}, third edition, Wiley 2024, chap.~6, pp.~99--125, with Exhibits~6.2, 6.4--6.7 by page image; and the authors' version of 5 January 2024 of \texttt{wide-boyd-markowitz-seventy}, secs.~1.1--1.4, 2.4, 3.2, 3.4, 4.3 and 5.6, pp.~3--10, 17, 19--22, 27--28 and 39--40, not the journal pagination; Exhibit~6.1 weights, pair count and Sharpe scale invariance computed). Stability: draft. Currency: formulation timeless; solver timings as of 2024.
\end{qfmodel}

\begin{qfmethod}[Signal-portfolio allocation: weighting simulated factor portfolios rather than positions, as Narang describes Grinold and Kahn's approach]\label{meth:signal-portfolio-allocation}
\qffield{Statement} Narang describes an optimiser ``directly aimed at building a portfolio of signals, whereas most optimizers try to size positions'', attributing it to Grinold and Kahn's \emph{Active Portfolio Management} \citep[chap.~6, sec.~6.2.2.4, pp.~117--118]{s7-narang}:
\begin{enumerate}
\item build one \emph{factor portfolio} per type of alpha forecast (momentum, value, growth), usually rule-based, ``very often equal-weighted or equal risk-weighted'';
\item simulate each historically to obtain a return series;
\item treat those series as the instruments of an optimiser, with the risk model, transaction cost model, portfolio size and risk target as inputs;
\item set each position as the sum, over factors, of its weight in the factor portfolio times that factor's signal times the factor portfolio's weight.
\end{enumerate}
His example: 100 equally weighted stocks, factor weights 60 and 40 percent, signals $+1$ and $-1$ for one company, give $1\%\times(+1)\times60\%+1\%\times(-1)\times40\%=+0.2\%$ (p.~118, page image; checked).

\qffield{Conditions} The number of factor portfolios is ``typically much more manageable, usually not more than about 20, corresponding to the number of individual factors in the alpha model'' (p.~117). That number is Narang's practitioner statement, with no source; it depends on how many signals a model has. The factor-portfolio returns are simulated, so the optimiser's inputs inherit the backtests' assumptions and their overfitting risk (\hyperref[pit:parameter-overfitting]{parameter overfitting}).

\qffield{Practice} Interpretation: the method trades a large position-level estimation problem for a small one over a handful of return series, ``an easier hurdle to clear in terms of the amount of data needed'' (p.~117). The position weights that result are linear in the factor weights, so position limits and costs at stock level are not enforced unless they are added back.

\qffield{Pitfalls} Whether this is Grinold and Kahn's own procedure, and in what form, was not checked: their book (\texttt{alpha-grinold}, a reading of the section overview) is unread. The legacy ``what is optimised is then a handful of factor portfolios'' and ``not more than 20'' repeat Narang's account as fact.

\qffield{See also} \hyperref[mod:mean-variance-portfolio-optimisation]{mean--variance optimisation}; \hyperref[meth:rule-based-position-sizing]{rule-based position sizing} (how each factor portfolio is weighted); \hyperref[meth:backtest-evaluation-diagnostics]{backtest evaluation diagnostics}.

\qffield{Further reading} Grinold and Kahn, \emph{Active Portfolio Management} (\texttt{alpha-grinold}), not read. The legacy file it replaces is kept verbatim in the section's folder of \nolinkurl{notes/artefacts/}, as\\ \nolinkurl{portfolio-construction-models.tex}.

\qffield{Index} factor portfolio; signal portfolio; Grinold and Kahn; alpha factor weighting; strategy allocation; simulated factor returns; position aggregation

\qffield{Status} Verification: partly checked 2026-10-15 (against the text layer of a local copy of \texttt{s7-narang}, third edition, Wiley 2024, chap.~6, pp.~117--118, with p.~118 by page image; the example computed; Grinold and Kahn not read). Stability: draft. Currency: practitioner account as of 2024.
\end{qfmethod}
